\documentclass{ctexart}
% ctexart   适合短文档：section、subsection、subsubsection

\usepackage{R3S-tex-doc}

% ---------- 自定义设置 ----------
% 目录分栏设置
% \multicoltoctrue      % 目录分两栏
% \multicoltocfalse     % 目录不分栏（默认）
\multicoltoctrue
% 颜色框自动编号设置
% \thmwithsectiontrue   % 颜色框随 section 编号（默认）
% \thmwithsectionfalse  % 颜色框全局编号
\thmwithsectionfalse
% 显示答案与解析设置
% \showanswerstrue      % 显示答案（默认）
% \showanswerfalse      % 不显示答案
% \showexplanationstrue   % 显示解析（默认）
% \showexplanationsfalse  % 不显示解析
\showanswerstrue
\showexplanationsfalse
% ----------------------------

\begin{document}
\title{}                    % 标题
\author{}                   % 作者
\date{}                     % 日期                  
\maketitle                  % 生成标题页（前面三项设置会体现出来）
\newpage                    % 新起一页避免标题页后接着目录内容
\tableofcontents            % 生成目录页（会自动包含章节、小节、子小节等信息）
\newpage                    % 新起一页避免目录后直接跟着内容

% ----- 书写内容解构示例 -----
% \section{章节1}
%   \subsection{小节1}
%       \subsubsection{子小节1}
%       \subsubsection{子小节1}
%   \subsection{小节2}
%       \subsubsection{子小节2}
% \section{章节2}
%   \subsection{小节3}
%       \subsubsection{子小节3}
%   \subsection{小节4}
%       \subsubsection{子小节4}
% \section{章节3}
%   \subsection{小节5}
%       \subsubsection{子小节5}
%   \subsection{小节6}
%       \subsubsection{子小节6}
% ----------------------------------------

\appendix                   % 附录部分

\end{document}